\section{Cadre théorique et hypothèses}\label{sec:cadre}

Le mécanisme que nous cherchons à identifier est l'arbitrage, par un actionnaire personne physique
qui contrôle la distribution de sa société, entre se verser un euro de bénéfice immédiatement
imposé et le laisser dans la société pour qu'elle l'investisse. Le cadre théorique est celui des
nouvelles vues de l'impôt sur les dividendes, dans lequel l'investissement financé par les bénéfices
retenus n'est pas affecté par le taux d'imposition des dividendes \citep{auerbach1979,king1977,chetty2005}.
Nous le présentons par un modèle à deux dates qui rend le mécanisme explicite et dont les hypothèses
testables descendent directement.

Un actionnaire détient un euro de bénéfice après impôt sur les sociétés. Il peut le distribuer
aujourd'hui, le dividende étant taxé au taux $\tau_d$, et placer le net au-dehors à un rendement $r$
lui-même imposé au taux $\tau_i$ ; ou bien le laisser dans la société, qui l'investit dans un projet
de rendement marginal $\rho$ et le distribue à la date suivante. La première voie procure, à la date
suivante, une consommation nette de $(1-\tau_d)\bigl(1 + r(1-\tau_i)\bigr)$. La seconde procure
$(1-\tau_d)(1+\rho)$. L'actionnaire est indifférent lorsque les deux voies s'égalisent,
\begin{equation}\label{eq:arbitrage}
  (1-\tau_d)\,(1+\rho) = (1-\tau_d)\,\bigl(1 + r(1-\tau_i)\bigr),
\end{equation}
où $\tau_d$ est le taux d'imposition des dividendes, $\tau_i$ celui des autres revenus du capital de
l'actionnaire, $r$ le rendement de son placement extérieur et $\rho$ le rendement marginal du projet
financé par la société. Le taux $\tau_d$ apparaît en facteur des deux membres et s'élimine, de sorte
que le rendement marginal qui déclenche l'investissement vaut
\begin{equation}\label{eq:neutralite}
  \rho^{\star} = r\,(1-\tau_i),
\end{equation}
indépendamment de $\tau_d$. C'est le résultat de neutralité des nouvelles vues : un euro de bénéfice
retenu supporte la taxe sur les dividendes qu'il soit distribué aujourd'hui ou demain, la taxe se
capitalise des deux côtés et ne déplace pas la marge d'investissement. Une variation de $\tau_d$ ne
change donc ni l'entrée de capital neuf, qui ne serait utile que pour financer un investissement
que $\rho^{\star}$ ne commande pas, ni le besoin de financement externe par la dette. Elle ne joue
que sur le calendrier et le niveau des versements. De là découlent les deux premières hypothèses.

\medskip
\noindent\textit{Hypothèse H1. Une baisse de la fiscalité des dividendes ne fait pas entrer de
capital neuf dans les fonds propres des sociétés.}
\smallskip

\noindent Le mécanisme est celui de l'équation~\eqref{eq:neutralite} : la réforme abaisse le coût de
sortie des fonds sans toucher le rendement $\rho^{\star}$ qui commande l'investissement, donc sans
créer de besoin de capital neuf. Le signe attendu sur l'entrée de capital est nul.

\medskip
\noindent\textit{Hypothèse H2. Une baisse de la fiscalité des dividendes ne provoque ni substitution
par la dette ni baisse de l'investissement.}
\smallskip

\noindent Le mécanisme est que, si rien ne quitte le financement interne pour un investissement
supplémentaire, les sociétés n'ont pas à remplacer les fonds distribués par de la dette. Les signes
attendus sur le levier et sur l'investissement sont nuls.

La troisième hypothèse est l'apport central et exige d'enrichir le modèle statique d'une dynamique
d'ajustement. \citet{bach2023} établissent que la hausse d'impôt de 2013 a accru la rétention des
bénéfices des sociétés contrôlées par leur actionnaire, un déplacement vers l'épargne interne plutôt
que vers l'investissement. Notons $R_t$ le stock de capital retenu dans la société à la date $t$ et
$R^{\star}(\tau_d)$ sa cible, croissante en $\tau_d$ puisqu'un impôt plus lourd sur la sortie incite
à conserver les fonds dans la société. Le stock s'ajuste vers sa cible selon
\begin{equation}\label{eq:ajustement}
  R_t - R_{t-1} = \lambda^{\pm}\,\bigl(R^{\star}(\tau_d) - R_{t-1}\bigr),
  \qquad \lambda^{+} \ge \lambda^{-},
\end{equation}
où $R_t$ est le capital retenu, $R^{\star}(\tau_d)$ sa cible fonction du taux d'imposition des
dividendes, et $\lambda^{\pm}$ la vitesse d'ajustement, qui prend la valeur $\lambda^{+}$ quand la
cible monte, après une hausse d'impôt qui pousse à retenir, et $\lambda^{-}$ quand elle baisse, après
une baisse d'impôt qui permet de relâcher. Les coûts d'ajustement et un motif de précaution rendent
la rétention plus rapide que la libération, soit $\lambda^{+} \ge \lambda^{-}$, ce qui engendre un
cliquet : à chocs d'ampleur égale et de sens opposé, le capital retenu réagit davantage à la hausse
qu'à la baisse. La même asymétrie se transmet à la distribution, miroir de la rétention.

\medskip
\noindent\textit{Hypothèse H3. La réponse de la distribution et de l'allocation du capital à la
fiscalité des dividendes est asymétrique : la baisse de 2018 ne défait qu'une fraction de l'effet de
la hausse de 2013, de sorte que le capital part plus facilement qu'il ne revient.}
\smallskip

\noindent Le mécanisme est celui de l'équation~\eqref{eq:ajustement} avec $\lambda^{+} > \lambda^{-}$.
Le signe attendu est une semi-élasticité de la distribution à la baisse inférieure, en valeur
absolue, à celle de la hausse, et une réallocation du capital plus faible au retour qu'à l'aller.

Ce même cadre oriente l'identification. L'arbitrage de l'équation~\eqref{eq:arbitrage} suppose un
actionnaire résident qui pilote la distribution, ce qui ne vaut que dans les sociétés privées
contrôlées. Dans une société cotée à l'actionnariat dispersé et international, aucun contribuable
résident ne fixe la politique de versement, qui obéit au signal de marché et non au taux marginal
$\tau_d$ d'un ménage. La réponse à la réforme ne peut donc vivre que dans les sociétés privées, ce
qui fournit le contrôle négatif des sociétés cotées.
